% Para fijar otra posición de la homotopía, cambie \tHom (0<tHom<1).
\begin{tikzpicture}[x=1cm,y=1cm,line width=.75pt,
                    line cap=round,line join=round,>=stealth]
  \def\L{2.2}       % Longitud común de los cilindros.
  \def\separacion{5}
  \def\tHom{.38}    % Valor de t representado.
  \def\alturaImagen{.40} % Distancia sobre el cilindro más ancho.

  \def\tapaVisible#1#2#3{%
    \draw (#1,0) ellipse [x radius=#2,y radius=#3];%
  }
  \def\tapaDerecha#1#2#3{%
    \draw[dashed] (#1,#3)
      arc[start angle=90,end angle=270,
          x radius=#2,y radius=#3];
    \draw (#1,#3)
      arc[start angle=90,end angle=-90,
          x radius=#2,y radius=#3];%
  }
  \def\tapaIzquierda#1#2#3#4{%
    \pgfmathsetmacro{\anguloOculto}{asin(#4/#3)}%
    \draw (#1,#3)
      arc[start angle=90,end angle={180-\anguloOculto},
          x radius=#2,y radius=#3];
    \draw[dashed] ({#1-#2*cos(\anguloOculto)},#4)
      arc[start angle={180-\anguloOculto},
          end angle={180+\anguloOculto},
          x radius=#2,y radius=#3];
    \draw ({#1-#2*cos(\anguloOculto)},{-#4})
      arc[start angle={180+\anguloOculto},end angle=450,
          x radius=#2,y radius=#3];%
  }
  \def\cuerpo#1#2{%
    \draw ({#1*\L},#2) -- ({(#1+1)*\L},#2)
          ({#1*\L},{-#2}) -- ({(#1+1)*\L},{-#2});%
  }
  \def\cuerpoAbierto#1#2{%
    \draw ({#1*\L},#2) -- ({(#1+.85)*\L},#2)
          ({#1*\L},{-#2}) -- ({(#1+.85)*\L},{-#2});
    \draw[dotted]
          ({(#1+.85)*\L},#2) -- ({(#1+1.2)*\L},#2)
          ({(#1+.85)*\L},{-#2}) -- ({(#1+1.2)*\L},{-#2});%
  }

  % Telescopio superior: X_n, Omega X_{n+1}, ...
  \cuerpo{0}{.38}
  \tapaVisible{0}{.09}{.38}
  \tapaDerecha{\L}{.09}{.38}
  \node[above=2pt] at (0,.38) {$X_n$};

  \cuerpo{1}{.72}
  \tapaIzquierda{\L}{.15}{.72}{.38}
  \tapaDerecha{2*\L}{.15}{.72}
  \node[above=2pt] at (\L,.72) {$\Omega X_{n+1}$};

  \cuerpo{2}{1.08}
  \tapaIzquierda{2*\L}{.22}{1.08}{.72}
  \tapaDerecha{3*\L}{.22}{1.08}
  \node[above=2pt] at ({2*\L},1.08) {$\Omega^2X_{n+2}$};

  \cuerpoAbierto{3}{1.50}
  \tapaIzquierda{3*\L}{.30}{1.50}{1.08}
  \node[above=2pt] at ({3*\L},1.50) {$\Omega^3X_{n+3}$};

  % Telescopio inferior: se omite X_n y se añade Omega^4 X_{n+4}.
  \begin{scope}[yshift=-\separacion cm]
    \cuerpo{0}{.72}
    \tapaVisible{0}{.15}{.72}
    \tapaDerecha{\L}{.15}{.72}
    \node[below=2pt] at (0,-.72) {$\Omega X_{n+1}$};

    \cuerpo{1}{1.08}
    \tapaIzquierda{\L}{.22}{1.08}{.72}
    \tapaDerecha{2*\L}{.22}{1.08}
    \node[below=2pt] at (\L,-1.08) {$\Omega^2X_{n+2}$};

    \cuerpo{2}{1.50}
    \tapaIzquierda{2*\L}{.30}{1.50}{1.08}
    \tapaDerecha{3*\L}{.30}{1.50}
    \node[below=2pt] at ({2*\L},-1.50) {$\Omega^3X_{n+3}$};

    \cuerpoAbierto{3}{1.86}
    \tapaIzquierda{3*\L}{.37}{1.86}{1.50}
    \node[below=2pt] at ({3*\L},-1.86) {$\Omega^4X_{n+4}$};

    % Segmentos de llegada por encima del cilindro más ancho
    % que interviene en cada imagen: forman una escalerita.
    % El primero representa [0,t] del primer cilindro inferior.
    \draw[rojoCIMAT,very thick]
      (0,{.72+\alturaImagen})
      -- ({\tHom*\L},{.72+\alturaImagen});
    \foreach \x in {0,\tHom}{
      \draw[rojoCIMAT] ({\x*\L},{.72+\alturaImagen-.08})
              -- ({\x*\L},{.72+\alturaImagen+.08});
    }
    \node[rojoCIMAT,below=-1.5pt,font=\scriptsize]
      at ({\tHom*\L},{.72+\alturaImagen-.08}) {$t$};

    % Los demás atraviesan una pegadura:
    % [t,1] del cilindro anterior unido a [0,t] del siguiente.
    \foreach \m/\radio in {1/1.08,2/1.50,3/1.86}{
      \pgfmathsetmacro{\ySegmento}{\radio+\alturaImagen}
      \draw[rojoCIMAT,very thick]
        ({(\m-1+\tHom)*\L},\ySegmento)
        -- ({(\m+\tHom)*\L},\ySegmento);
      \foreach \extremo in {0,1}{
        \draw[rojoCIMAT]
          ({(\m-1+\tHom+\extremo)*\L},{\ySegmento-.08})
          -- ({(\m-1+\tHom+\extremo)*\L},{\ySegmento+.08});
      }
      \node[rojoCIMAT,below=-1.5pt,font=\scriptsize]
        at ({(\m+\tHom)*\L},{\ySegmento-.08}) {$t$};
    }
  \end{scope}

  % Una flecha por cilindro superior; cada flecha termina
  % exactamente sobre su segmento rojo de llegada.
  \foreach \m/\radioSuperior/\radioInferior in
    {0/.38/.72,1/.72/1.08,2/1.08/1.50,3/1.50/1.86}{
    \pgfmathsetmacro{\xImagen}
      {(\m==0 ? .5*\tHom : \m-.5+\tHom)*\L}
    \pgfmathsetmacro{\yImagen}
      {-\separacion+\radioInferior+\alturaImagen}
    \pgfmathsetmacro{\yControl}{(-\radioSuperior+\yImagen)/2}
    \draw[rojoCIMAT,thick,->]
      ({(\m+.5)*\L},{-\radioSuperior})
      .. controls ({(\m+.5)*\L},\yControl)
               and (\xImagen,\yControl)
      .. (\xImagen,\yImagen);
  }
  \node[rojoCIMAT] at (8.65,-2.35) {$H_t$};
\end{tikzpicture}